\documentclass[11pt,a4paper]{article}
\usepackage[margin=2.4cm]{geometry}
\usepackage[T1]{fontenc}
\usepackage[utf8]{inputenc}
\usepackage{listings,xcolor,hyperref,booktabs,enumitem,titlesec,parskip}
\usepackage[expansion=false]{microtype}

% ---- colour palette --------------------------------------------------------
\definecolor{leafcyan}{RGB}{0,140,140}
\definecolor{leafgray}{RGB}{240,240,240}
\definecolor{leafdark}{RGB}{40,40,40}
\definecolor{leafcomment}{RGB}{100,120,100}

\hypersetup{colorlinks=true,linkcolor=leafcyan,urlcolor=leafcyan,citecolor=leafcyan}

% ---- listings --------------------------------------------------------------
\lstdefinestyle{leaf}{
  basicstyle=\ttfamily\footnotesize,
  breaklines=true,
  backgroundcolor=\color{leafgray},
  frame=single,
  rulecolor=\color{gray!40},
  commentstyle=\color{leafcomment}\itshape,
  keywordstyle=\color{leafcyan}\bfseries,
  stringstyle=\color{orange!70!black},
  xleftmargin=8pt,xrightmargin=8pt,
  aboveskip=7pt,belowskip=5pt,
  showstringspaces=false,
}
\lstset{style=leaf}

% ---- section formatting ----------------------------------------------------
\titleformat{\section}{\large\bfseries\color{leafcyan}}{}{0pt}{}[\titlerule]
\titleformat{\subsection}{\normalsize\bfseries\color{leafdark}}{}{0pt}{}

% ---- doc meta --------------------------------------------------------------
\title{%
  \textcolor{leafcyan}{\rule{\linewidth}{1.5pt}}\\[8pt]
  \textbf{LeafOS Optional Node System}\\[4pt]
  \large Advanced Learning Guide\\[2pt]
  \normalsize\texttt{v0.5.0}\\[6pt]
  \textcolor{leafcyan}{\rule{\linewidth}{0.8pt}}
}
\author{}
\date{}

\begin{document}
\maketitle\thispagestyle{empty}

\begin{center}
\fcolorbox{leafcyan}{leafgray}{\parbox{0.86\linewidth}{
\textbf{Scope note.} This guide covers the optional SSH node subsystem. The
current project-agent path is \texttt{leafctl live PATH}, backed by the
version-2 typed inlet, resident supervisor, CPU-authoritative validation, and
Python/native TUI. Read \texttt{DOCUMENTATION\_MAP.md},
\texttt{ARCHITECTURE.md}, and \texttt{RESIDENT\_STACK\_USAGE.md} first for
local project automation.
}}
\end{center}

\tableofcontents
\newpage

%% ===========================================================================
\section{Architecture Overview}

\subsection{One binary, two roles}

\texttt{leaf} is a single executable that branches on the command group invoked.
The receiver engine (\path{core/node/leaf_node.py}) runs locally on every node.
The distributor layer (\path{core/remote/remote.sh}) drives remote nodes by
SSH-ing into them and invoking the same binary.

\begin{lstlisting}[language={}]
MASTER / distributor
  bin/leafctl --source--> core/remote/remote.sh
								 |
						  LAN / WAN SSH
								 |
WORKER / receiver
  ssh -> bin/leafctl --source--> core/node/node.sh
									   |
							   python3 leaf_node.py
									   |
  ~/.leaf/  inbox/  queued/  running/  done/  failed/
			logs/TASK_ID.jsonl
			results/TASK_ID/result.json
\end{lstlisting}

The master never contacts an HTTP API. It opens one SSH connection per operation
and reads files the receiver wrote. That is the entire protocol.

\subsection{Source chain}

\begin{lstlisting}[language=bash]
bin/leafctl
  source core/brand/brand.sh
  source core/log/log.sh
  source core/remote/remote.sh      # distributor layer
	source core/node/node.sh        # receiver wrapper
	  resolves python3
	  exports LEAF_NODE_PY path
\end{lstlisting}

\texttt{leaf\_node.py} is loaded fresh per invocation. A node doing no work is
a directory tree and an SSH daemon. No persistent process.

\subsection{Workspace resolution}

\begin{lstlisting}[language={}]
leaf_home() resolution order:
  1. --home FLAG
  2. $LEAF_HOME environment variable
  3. ~/.leaf  (default)
\end{lstlisting}

%% ===========================================================================
\section{State Machine}

\subsection{Task lifecycle}

\begin{lstlisting}[language={}]
							+---- cancel flag ---+
							|                    v
inbox/ -> queued/ -> running/ -> done/      failed/
						 |
						 +-> failed/   (non-zero exit)
\end{lstlisting}

State transitions are atomic file moves (\texttt{shutil.move}). A task occupies
exactly one directory at all times. Crashes abandoning \texttt{running/} are
detectable by \texttt{leaf task list --state running}.

\subsection{Cooperative cancellation}

The distributor creates \texttt{running/TASK\_ID.cancel}. The runner polls for
it. A runner that ignores the flag finishes normally --- there is no
\texttt{SIGKILL} without warning.

\begin{lstlisting}[language=bash]
leaf task cancel gpu1 task_20260621_142203
# internally (over SSH):
# touch ~/.leaf/running/task_20260621_142203.cancel
\end{lstlisting}

In Python, check with:
\begin{lstlisting}[language=python]
if os.path.exists(_cancel_flag(home, task["task_id"])):
	return 130, "cancelled", {}
\end{lstlisting}

%% ===========================================================================
\section{Writing a Custom Runner}

A runner is a Python function:

\begin{lstlisting}[language=python]
def run_myrunner(home: str, task: dict, log) -> tuple:
	"""
	home  -- absolute path to $LEAF_HOME
	task  -- parsed task JSON descriptor
	log   -- TaskLog; log.event("stdout", text="...")
	Returns (exit_code: int, status: str, output: dict)
	"""
\end{lstlisting}

\subsection{Minimal example: sleep runner}

\begin{lstlisting}[language=python]
import time, os

def run_sleep(home, task, log):
	secs = float(task.get("payload", {}).get("seconds", 5))
	log.event("stdout", text="sleeping {}s".format(secs))
	time.sleep(secs)
	return 0, "ok", {"slept": secs}
\end{lstlisting}

Register in \texttt{RUNNERS} in \texttt{leaf\_node.py}:
\begin{lstlisting}[language=python]
RUNNERS = {
	"echo": run_echo, "shell": run_shell,
	"hardware": run_hardware, "llama": run_llama,
	"sleep": run_sleep,   # add here
}
\end{lstlisting}

\subsection{Streaming stdout with real-time cancellation}

\begin{lstlisting}[language=python]
def run_myrunner(home, task, log):
	import subprocess
	proc = subprocess.Popen(
		["long-running-cmd"],
		stdout=subprocess.PIPE, stderr=subprocess.STDOUT,
		universal_newlines=True,
	)
	for line in iter(proc.stdout.readline, ""):
		log.event("stdout", text=line.rstrip("\n"))
		if os.path.exists(_cancel_flag(home, task["task_id"])):
			proc.terminate()
			return 130, "cancelled", {}
	code = proc.wait()
	return code, "ok" if code == 0 else "error", {"exit_code": code}
\end{lstlisting}

The master sees each line in real time via \texttt{leaf task watch} (\texttt{tail -f} over SSH).

%% ===========================================================================
\section{The Three Contracts}

These schemas must stay stable. Internal rewrites and new runners are safe as
long as these hold.

\subsection{Task file schema}

\begin{lstlisting}[language={}]
{
  "task_id":  "task_20260621_142203",  // stamped by distributor
  "kind":     "echo",                  // selects the runner
  "priority": 5,                       // 1 (low) - 10 (high)
  "payload":  { "message": "hello" }  // runner-defined
}
\end{lstlisting}

\subsection{State directory names (frozen)}

\begin{center}
\texttt{inbox} \quad \texttt{queued} \quad \texttt{running} \quad
\texttt{done} \quad \texttt{failed}
\end{center}

Do not rename. \texttt{\_find\_task()} iterates \texttt{STATE\_DIRS} in order.

\subsection{JSONL log event schema}

\begin{tabular}{@{}llc@{}}
\toprule
field & type & always present \\
\midrule
\texttt{time}      & ISO-8601 string & yes \\
\texttt{event}     & string          & yes \\
\texttt{task\_id}  & string          & yes \\
\texttt{runner}    & string          & \texttt{started} only \\
\texttt{text}      & string          & \texttt{stdout}/\texttt{stderr} \\
\texttt{exit\_code}& int             & terminal events \\
\bottomrule
\end{tabular}

\medskip
Standard sequence:
\texttt{queued \(\to\) started \(\to\) stdout* \(\to\) done\,|\,failed\,|\,cancelled}

%% ===========================================================================
\section{Transport Layer Internals}

\subsection{Registry (nodes.json)}

\begin{lstlisting}[language={}]
{
  "gpu1": {
	"target":    "user@192.168.1.44",
	"path":      "~/.leaf",
	"transport": "ssh"
  }
}
\end{lstlisting}

\texttt{\_node\_resolve NAME} in \texttt{remote.sh} reads this via
\texttt{registry-get} (Python, no \texttt{jq} dependency), exporting
\texttt{\_RT}, \texttt{\_RTGT}, \texttt{\_RPATH}.

\subsection{Transport dispatch in remote.sh}

\begin{lstlisting}[language=bash]
case "$_RT" in
	local) _local_engine task-start "$task_id" ;;
	ssh)   _ssh "$_RTGT" "leaf task start $task_id" ;;
esac
\end{lstlisting}

\texttt{\_local\_engine} is \texttt{LEAF\_HOME="\$\_RPATH" python3 leaf\_node.py}.
This is why \texttt{tests/nodes.sh} needs no network.

\subsection{SSH options}

\begin{lstlisting}[language=bash]
export LEAF_SSH_OPTS="-i ~/.ssh/leaf_node_key -p 2222"
\end{lstlisting}

Expanded by \texttt{\_ssh()} and \texttt{\_scp()} before all other arguments.

%% ===========================================================================
\section{Multi-Node Patterns}

\subsection{Fan-out: submit to multiple nodes}

\begin{lstlisting}[language=bash]
for node in gpu1 gpu2 cpu1; do
	TASK_ID=$(leaf task submit "$node" examples/echo.task.json | tail -n1)
	echo "$node $TASK_ID"
done
\end{lstlisting}

\subsection{Collect results from all workers}

\begin{lstlisting}[language=bash]
mkdir -p ./run_collection
for node in gpu1 gpu2; do
	leaf task pull "$node" "$TASK_ID" "./run_collection/$node"
done
\end{lstlisting}

\subsection{Choose the idle node}

\begin{lstlisting}[language=bash]
pick_idle() {
	leaf nodes list --json | python3 -c "
import json, sys, subprocess
nodes = json.load(sys.stdin)
for name, cfg in nodes.items():
	r = subprocess.run(['leaf','nodes','ping',name],
					   capture_output=True, text=True)
	d = json.loads(r.stdout) if r.returncode == 0 else {}
	if d.get('tasks_running', 1) == 0:
		print(name); break
"
}
TARGET=$(pick_idle)
leaf task submit "$TARGET" examples/shell.task.json
\end{lstlisting}

%% ===========================================================================
\section{Security Hardening}

\subsection{Dedicated SSH key}

\begin{lstlisting}[language=bash]
ssh-keygen -t ed25519 -C "leaf-cluster" -f ~/.ssh/leaf_node_key
ssh-copy-id -i ~/.ssh/leaf_node_key.pub user@192.168.1.44
export LEAF_SSH_OPTS="-i ~/.ssh/leaf_node_key"
\end{lstlisting}

\subsection{Restrict key to leaf commands only}

In \texttt{\~{}/.ssh/authorized\_keys} on the worker:

\begin{lstlisting}[language={}]
command="leaf task start ${SSH_ORIGINAL_COMMAND#* }",
no-port-forwarding,no-X11-forwarding,no-agent-forwarding,no-pty
ssh-ed25519 AAAA... leaf-cluster
\end{lstlisting}

\subsection{Firewall (ufw example)}

\begin{lstlisting}[language=bash]
ufw allow from 192.168.1.0/24 to any port 22
ufw deny 22
\end{lstlisting}

No other port needs to be open. No HTTP listener, no agent socket, no
exposed API --- only SSH.

\subsection{Non-root dedicated user}

\begin{lstlisting}[language=bash]
useradd -m -s /bin/bash leafrunner
su - leafrunner -c "leaf node init --node-id worker1"
\end{lstlisting}

%% ===========================================================================
\section{Debugging and Observability}

\subsection{Read JSONL logs directly}

\begin{lstlisting}[language=bash]
# all events
cat ~/.leaf/logs/TASK_ID.jsonl | python3 -m json.tool

# stdout lines only
grep '"event":"stdout"' ~/.leaf/logs/TASK_ID.jsonl \
  | python3 -c "
import sys, json
[print(json.loads(l)['text']) for l in sys.stdin]"
\end{lstlisting}

\subsection{Orphaned running tasks}

\begin{lstlisting}[language=bash]
leaf task list --state running --json
leaf task cancel TASK_ID   # drops cancel flag; runner exits on next poll
\end{lstlisting}

\subsection{LEAF\_HOME isolation for debugging}

\begin{lstlisting}[language=bash]
LEAF_HOME=/tmp/debug_node leaf node init --node-id debug
LEAF_HOME=/tmp/debug_node leaf task start TASK_ID
\end{lstlisting}

\subsection{SSH connectivity check}

\begin{lstlisting}[language=bash]
ssh $LEAF_SSH_OPTS user@192.168.1.44 "leaf node status --json"
ssh -v -o ConnectTimeout=5 user@192.168.1.44 "echo ok"
\end{lstlisting}

%% ===========================================================================
\section{Extension Roadmap}

\begin{tabular}{@{}ll@{}}
\toprule
Milestone & Planned change \\
\midrule
0.6.0 & \texttt{llama.cpp} runner; node selection by GPU presence \\
0.7.0 & Priority queue; max-N concurrent tasks; queued drains by priority \\
0.8.0 & \texttt{rsync} replaces \texttt{scp} for large/partial payloads \\
0.9.0 & Restricted SSH command wrapper + \texttt{authorized\_keys} templates \\
1.0.0 & Optional read-only web viewer (\texttt{leaf node serve}), log streaming only \\
\bottomrule
\end{tabular}

\medskip
The three contracts in Section~4 are frozen from this milestone forward.

%% ===========================================================================
\section{Full Command Reference}

\subsection*{Receiver (runs on the worker)}

\begin{lstlisting}[language={}]
leaf node init   [--node-id ID] [--role ROLE] [--json]
leaf node status [--json]
leaf node hardware [--json]
leaf task start   TASK_ID
leaf task list    [--state STATE] [--json]
leaf task logs    TASK_ID [--follow]
leaf task result  TASK_ID [--json]
leaf task cancel  TASK_ID
\end{lstlisting}

\subsection*{Distributor (runs on the master)}

\begin{lstlisting}[language={}]
leaf nodes add   NAME USER@HOST [--transport ssh|local] [--path DIR]
leaf nodes list  [--json]
leaf nodes ping  NAME
leaf task submit NAME TASK_FILE
leaf task watch  NAME TASK_ID
leaf task pull   NAME TASK_ID [DEST]
leaf task cancel NAME TASK_ID
\end{lstlisting}

\subsection*{Key files}

\begin{tabular}{@{}ll@{}}
\toprule
File & Purpose \\
\midrule
\texttt{core/node/leaf\_node.py}  & Receiver engine (stdlib only) \\
\texttt{core/node/node.sh}        & Brand-aware receiver wrapper \\
\texttt{core/remote/remote.sh}    & Distributor + transport layer \\
\texttt{bin/leafctl}              & CLI dispatch \\
\texttt{examples/*.task.json}     & Example task descriptors \\
\texttt{tests/nodes.sh}           & End-to-end local-transport test suite \\
\texttt{docs/NODES.md}            & Architecture reference \\
\texttt{docs/QUICKSTART.md/.tex}  & Quick start \\
\bottomrule
\end{tabular}

\end{document}
